In this section, we discuss perfect utility maximizers with bounded knowledge and show performance guarantees for such agents. In \Cref{sec:combining}, we combine these results, show how to relax the assumption of perfect optimization and, most importantly, show how the performance of a bounded-rational agent differs between the cases with and without self-modification.

In \Cref{lem:optimization_bound}, we show that if the agent's estimate of the expected discounted utility is at most $\epsilon$ away from the true value, the agent will be a $2\epsilon$-optimizer. In \crefrange{sec:utility}{sec:gammas}, we show bounds on how inaccurate the agent's estimate of expected discounted utility can be, thus proving bounds on optimization. In \Cref{sec:gammas}, we proceed differently: we formulate the worst case as a solution of an optimization problem which we then solve analytically.

\begin{lemma} \label{lem:optimization_bound}
Let A be a (possibly self-modifying) perfect expected utility maximizer with knowledge $\kappa = (u,\rho, \gamma)$. Let $\kappa^* = (u^*,\rho^*, \gamma^*)$ be the correct knowledge. Assume that
\begin{equation} \label{eq:bound_assumption}
|V^\pi_\kappa(\mae_{<t}) - V^\pi_{\kappa^*}(\mae_{<t})| \leq \epsilon
\end{equation}
for all policies $\pi$ and histories $\mae_{<t}$. Then, A is a $2\epsilon$-optimizer with respect to $\kappa^*$.
\end{lemma}

